\section{Inclusion operators and normalized compatible fields}
\label{organized:inclusion-fields}
Write $\alpha$ for the kernel partition distribution and $\beta$ for its
range distribution. A subset is split by a partition if its vertices lie
in different blocks. For $1\le\ell\le r$, define
\[
 u_\ell(I)=\Pr_\alpha\{I\text{ is split}\},\qquad
 \eta_\ell(I)=\binom r\ell^{-1}\Pr_\beta\{I\subseteq\mathcal R\}.
\]
Thus $\eta_\ell$ is a probability row, with $\eta_r=\beta$. For adjacent
levels the normalized range marginal is
\[
 \eta_\ell(I)=\frac1{\ell+1}
 \sum_{j\notin I}\eta_{\ell+1}(I\cup\{j\}).
\]
Indeed every $(\ell+1)$-set contributes to exactly $\ell+1$ of its
$\ell$-subsets. Equivalently, count the $r-\ell$ extensions inside a
fixed range and use the binomial coefficients.

Equipartition gives the exact right-field normalization
\[
 u_\ell(I)=\frac r{r-\ell}
 \sum_{j\notin I}\pi_j u_{\ell+1}(I\cup\{j\}),\qquad\ell<r.
\]
For a fixed kernel partition, an already split $I$ occupies $\ell$ blocks;
the union of the unoccupied blocks has stationary mass $(r-\ell)/r$.
If $I$ is not split, no extension is split. Average this identity over
$\alpha$ to obtain the formula.

The incidence matrix $E_{\ell+1,\ell}(J,I)=\mathbf1_{I\subset J}$
expresses the first recurrence in operator form. The common-letter lift
maps each kernel range injectively to another range, so its top stationary
range law descends to stationary $\eta_\ell$ at every level.
The right field is $u_\ell=\Omega_\ell\mathbf1$, with the Abel projection
$\Omega_\ell$. One must not infer that every intermediate projection has
rank one, or identify a normalized row $\mathbf1^{\mathsf T}\Omega_\ell$
with the chosen top-range marginal without an argument.
The full inclusion-compatibility identities and their proofs have their
primary source in \texttt{kernelHierarchy.tex}; this section fixes the
normalizations used in the organized edition.
